\section{Sequence Diagram}

\usetikzlibrary{positioning,arrows.meta,calc}

Phần này mô tả tuần tự các luồng hoạt động chính của hệ thống thông qua các sơ đồ trình tự (sequence diagram).

% =========================================================
\subsection{Luồng đăng nhập và kết nối P2P}
% =========================================================

Hình~\ref{fig:seq-p2p} mô tả toàn bộ quá trình từ khi Client A đăng nhập đến khi thiết lập kết nối P2P với Client B và trao đổi tin nhắn.

\begin{figure}[H]
\centering
\begin{tikzpicture}[
    font=\scriptsize,
    actor/.style={draw, fill=gray!15, minimum width=2.3cm, minimum height=0.6cm, align=center},
    msg/.style={-{Stealth[length=4pt]}, thick},
    notelbl/.style={draw, fill=yellow!20, rounded corners=2pt, inner sep=3pt, font=\tiny, align=left}
]

\def\xA{0}
\def\xS{5.5}
\def\xB{11}

\node[actor] at (\xA, 10) {Client A};
\node[actor] at (\xS, 10) {Server};
\node[actor] at (\xB, 10) {Client B};

\draw[dashed, gray] (\xA, 9.7) -- (\xA, -0.3);
\draw[dashed, gray] (\xS, 9.7) -- (\xS, -0.3);
\draw[dashed, gray] (\xB, 9.7) -- (\xB, -0.3);

%% --- Login A ---
\draw[msg] (\xA,9.3)--(\xS,9.3) node[midway,above,font=\tiny]{HELLO};
\draw[msg] (\xS,9.0)--(\xA,9.0) node[midway,above,font=\tiny]{HELLO\_ACK};
\draw[msg] (\xA,8.7)--(\xS,8.7) node[midway,above,font=\tiny]{AUTH\_REQUEST (alice)};
\draw[msg] (\xS,8.4)--(\xA,8.4) node[midway,above,font=\tiny]{AUTH\_RESPONSE (OK)};

%% --- Login B ---
\draw[msg] (\xB,9.3)--(\xS,9.3) node[midway,above,font=\tiny]{HELLO};
\draw[msg] (\xS,9.0)--(\xB,9.0) node[midway,above,font=\tiny]{HELLO\_ACK};
\draw[msg] (\xB,8.7)--(\xS,8.7) node[midway,above,font=\tiny]{AUTH\_REQUEST (bob)};
\draw[msg] (\xS,8.4)--(\xB,8.4) node[midway,above,font=\tiny]{AUTH\_RESPONSE (OK)};

%% --- Peer Lookup ---
\draw[msg] (\xA,7.9)--(\xS,7.9) node[midway,above,font=\tiny]{PEER\_LOOKUP\_REQUEST (bob)};
\draw[msg] (\xS,7.6)--(\xA,7.6) node[midway,above,font=\tiny]{PEER\_LOOKUP\_RESPONSE (IP:peerPort)};

%% --- P2P Handshake ---
\node[notelbl] at (3.0,7.1) {A mở TCP trực tiếp\\đến peerPort của B};
\draw[msg, red!70!black] (\xA,6.7)--(\xB,6.7) node[midway,above,font=\tiny]{PEER\_HELLO (alice)};
\draw[msg, red!70!black] (\xB,6.4)--(\xA,6.4) node[midway,above,font=\tiny]{PEER\_HELLO\_ACK};

%% --- Chat ---
\draw[msg, red!70!black] (\xA,5.9)--(\xB,5.9) node[midway,above,font=\tiny]{CHAT\_MESSAGE ``Xin chào Bob!''};
\draw[msg, red!70!black] (\xB,5.6)--(\xA,5.6) node[midway,above,font=\tiny]{CHAT\_MESSAGE ``Chào Alice!''};
\node[notelbl] at (5.5,5.1) {Tin nhắn P2P trực tiếp,\\không qua server};

\end{tikzpicture}
\caption{Sequence diagram: Đăng nhập và nhắn tin P2P trực tiếp}
\label{fig:seq-p2p}
\end{figure}

% =========================================================
\subsection{Luồng truyền file P2P}
% =========================================================

Hình~\ref{fig:seq-file-p2p} mô tả chi tiết quá trình truyền file trực tiếp giữa hai client sau khi đã thiết lập kết nối P2P.
Toàn bộ luồng dữ liệu không đi qua server.

\begin{figure}[H]
\centering
\begin{tikzpicture}[
    font=\scriptsize,
    actor/.style={draw, fill=gray!15, minimum width=2.6cm, minimum height=0.6cm, align=center},
    msg/.style={-{Stealth[length=4pt]}, thick},
    notelbl/.style={draw, fill=yellow!20, rounded corners=2pt, inner sep=3pt, font=\tiny, align=left}
]

\def\xA{0}
\def\xB{10}

\node[actor] at (\xA, 10) {Client A (sender)};
\node[actor] at (\xB, 10) {Client B (PeerListener)};

\draw[dashed,gray] (\xA,9.7)--(\xA,-0.3);
\draw[dashed,gray] (\xB,9.7)--(\xB,-0.3);

%% Existing P2P connection: A sends FILE_OFFER
\draw[msg] (\xA,9.3)--(\xB,9.3) node[midway,above,font=\tiny]{FILE\_OFFER (tên, kích thước, transferId)};
\node[notelbl] at (\xB+1.6,8.85) {B lưu offer \#N\\(transferId, sender)};
\draw[msg] (\xB,8.6)--(\xA,8.6) node[midway,above,font=\tiny]{FILE\_OFFER\_ACK};

\node[notelbl] at (\xB+1.6,8.1) {Người dùng gõ \lstinline{/accept N}};
\draw[msg] (\xB,7.7)--(\xA,7.7) node[midway,above,font=\tiny]{FILE\_DECISION (ACCEPT, transferId)};

%% New TCP connection for data
\node[notelbl] at (5.0,7.2) {A mở kết nối TCP mới đến peerPort của B};
\draw[msg] (\xA,6.9)--(\xB,6.9) node[midway,above,font=\tiny]{FILE\_OFFER (requestId=0) — header luồng data};

%% Chunks
\draw[msg] (\xA,6.5)--(\xB,6.5) node[midway,above,font=\tiny]{FILE\_CHUNK \#1};
\draw[msg] (\xA,6.2)--(\xB,6.2) node[midway,above,font=\tiny]{FILE\_CHUNK \#2 \ldots};
\draw[msg] (\xA,5.9)--(\xB,5.9) node[midway,above,font=\tiny]{FILE\_CHUNK \#N};

%% End and result
\draw[msg] (\xA,5.5)--(\xB,5.5) node[midway,above,font=\tiny]{FILE\_END (SHA-256 checksum)};
\draw[msg] (\xB,5.2)--(\xA,5.2) node[midway,above,font=\tiny]{FILE\_RESULT (OK / CHECKSUM\_MISMATCH)};
\node[notelbl] at (5.0,4.7) {B kiểm tra SHA-256, atomic move\\từ file tạm sang đích};

\end{tikzpicture}
\caption{Sequence diagram: Truyền file P2P trực tiếp}
\label{fig:seq-file-p2p}
\end{figure}

% =========================================================
\subsection{Luồng broadcast file qua server}
% =========================================================

Hình~\ref{fig:seq-broadcast} mô tả luồng broadcast file trong cộng đồng: Client A upload file lên server,
server lưu trữ và phân phối đề nghị tải xuống đến các thành viên cùng cộng đồng.
Mỗi client nhận tải file xuống qua một kết nối TCP riêng biệt đến cổng 8080.

\begin{figure}[H]
\centering
\begin{tikzpicture}[
    font=\scriptsize,
    actor/.style={draw, fill=gray!15, minimum width=2.3cm, minimum height=0.6cm, align=center},
    msg/.style={-{Stealth[length=4pt]}, thick},
    notelbl/.style={draw, fill=yellow!20, rounded corners=2pt, inner sep=3pt, font=\tiny, align=left}
]

\def\xA{0}
\def\xS{5.5}
\def\xB{11}

\node[actor] at (\xA, 11.5) {Client A (sender)};
\node[actor] at (\xS, 11.5) {Server};
\node[actor] at (\xB, 11.5) {Client B/C (receiver)};

\draw[dashed,gray] (\xA,11.2)--(\xA,-0.3);
\draw[dashed,gray] (\xS,11.2)--(\xS,-0.3);
\draw[dashed,gray] (\xB,11.2)--(\xB,-0.3);

%% Upload request
\draw[msg] (\xA,10.8)--(\xS,10.8)
    node[midway,above,font=\tiny]{FILE\_BROADCAST\_UPLOAD\_REQUEST (meta, communityId)};
\draw[msg] (\xS,10.5)--(\xA,10.5)
    node[midway,above,font=\tiny]{FILE\_BROADCAST\_UPLOAD\_RESPONSE (uploadToken)};

%% New TCP upload
\node[notelbl] at (2.5,10.1) {A mở TCP mới đến port 8080};
\draw[msg] (\xA,9.8)--(\xS,9.8)
    node[midway,above,font=\tiny]{FILE\_BROADCAST\_UPLOAD\_START (uploadToken)};
\draw[msg] (\xA,9.5)--(\xS,9.5)
    node[midway,above,font=\tiny]{FILE\_CHUNK \#1 \ldots \#N};
\draw[msg] (\xA,9.2)--(\xS,9.2)
    node[midway,above,font=\tiny]{FILE\_END (SHA-256)};
\draw[msg] (\xS,8.9)--(\xA,8.9)
    node[midway,above,font=\tiny]{FILE\_RESULT (OK)};

\node[notelbl] at (5.5,8.4) {Server lưu file vào\\FileStorageService};

%% Server offers to community members
\draw[msg] (\xS,8.0)--(\xB,8.0)
    node[midway,above,font=\tiny]{FILE\_BROADCAST\_OFFER (fileId, tên, kích thước)};

%% Recipients accept
\node[notelbl] at (\xB+1.6,7.5) {Người dùng gõ \lstinline{/accept N}};
\draw[msg] (\xB,7.2)--(\xS,7.2)
    node[midway,above,font=\tiny]{FILE\_BROADCAST\_DECISION (ACCEPT, offerId)};
\draw[msg] (\xS,6.9)--(\xB,6.9)
    node[midway,above,font=\tiny]{FILE\_BROADCAST\_DOWNLOAD\_RESPONSE (downloadToken)};

%% New TCP download
\node[notelbl] at (8.5,6.5) {B/C mở TCP mới đến port 8080};
\draw[msg] (\xB,6.2)--(\xS,6.2)
    node[midway,above,font=\tiny]{FILE\_BROADCAST\_DOWNLOAD\_START (downloadToken)};
\draw[msg] (\xS,5.9)--(\xB,5.9)
    node[midway,above,font=\tiny]{FILE\_CHUNK \#1 \ldots \#N (FileSender stream)};
\draw[msg] (\xS,5.6)--(\xB,5.6)
    node[midway,above,font=\tiny]{FILE\_END (SHA-256)};
\draw[msg] (\xB,5.3)--(\xS,5.3)
    node[midway,above,font=\tiny]{FILE\_RESULT (OK)};

\node[notelbl] at (8.5,4.8) {B/C kiểm tra SHA-256\\atomic move sang đích};

\end{tikzpicture}
\caption{Sequence diagram: Broadcast file trong cộng đồng qua server trung tâm}
\label{fig:seq-broadcast}
\end{figure}
